% ---------------------------------------------------------------------------
% Section: Phase 14 -- the engine's knobs, the model switch, and a cascade in the harness
% ---------------------------------------------------------------------------
\section{Phase 14: Engine Knobs, the Model Switch, and a Cascade in the Harness}
\label{sec:phase14}

Phase~14 (1 November 2026) asked three questions of the system as it stood after Phase~13: what the inference engine's own
settings can still give on the target machine, how fast the request moves from one model to the other (and the service from
one backend to another), and what of Codex's agent core the harness still lacks. The research note is
\texttt{docs/research/engine\_switching\_harness.md} and the decision is \texttt{docs/decisions/008-engine-switching-harness.md}.
The method is the one the project has used since Phase~7: measure first on the reference machine, interleave the arms, and
keep only what the numbers justify. This phase ran while another workload of the owner's shared the four cores (load average
8--67 against about 2 when idle), which decides what can be concluded from the engine measurements (Section~\ref{sec:p14-engine}).

\subsection{Literature and what transfers}

Phases 7--13 covered what changes \emph{what} the engine computes: speculative decoding \cite{leviathan2023speculative} in
three forms, weight and cache formats, batching across slots \cite{yu2022orca}, and prefix caching \cite{kwon2023pagedattention,
zheng2024sglang}. Three directions remained. (i)~\emph{How} the engine computes: flash attention on a CPU, the micro-batch
size, the thread count, and the busy-wait level at which ggml's worker threads spin while waiting for each other, which matters
when two servers' eight threads share four cores. (ii)~The switch between models: with both models resident (1.2 and 2.8\,GB of
32), a switch is routing plus a warm prefix, not a load, so the checkpoint-loading systems built for GPU memory pressure
\cite{fu2024serverlessllm} do not apply; what applies is making the prefix warm before the switch. (iii)~A \emph{cascade}
\cite{chen2023frugalgpt,yue2024cascades} inside the harness: the quick model first, the thorough one only when needed, with the
exact verifier as the cascade's scorer, which is the gate's logic applied to the agent. Ternary models and table-lookup kernels
\cite{ma2025bitnet,wei2024tmac} would change the model, not the engine, and no Qwen3 variant exists; they were not pursued.

\subsection{The engine's knobs}
\label{sec:p14-engine}

\texttt{scripts/engine\_tuning.py} runs \texttt{llama-bench} on the production GGUF files under each flag: flash attention on and
off, micro-batches of 256 and 1,024 tokens, three threads, no polling (\texttt{--poll 0}), locked memory; prompt processing of
512 tokens and decode of 64 tokens at a depth of 1,024 (the depth of a repair prompt). A first, block-ordered run was discarded:
the owner's other workload started mid-way and turned single blocks by a factor of 0.4. The second run cycles the arms in three
rounds, one repetition each, so a burst of load hits every arm, and records the host load before every row.
Table~\ref{tab:p14-flags} gives the medians. The verdict is \emph{inconclusive, flags unchanged}: the per-round ratios of a
flag against the same round's baseline swing from 0.07$\times$ to 3.8$\times$ (System~2 decode under flash attention: 0.15, 0.07,
1.04), which is the external load, not the flag, and the baseline itself sits 2--4 times below its idle values (31 tokens/s
prefill and 5--6 tokens/s decode for System~2, Phase~13). No flag is adopted on this evidence. What remains is the method, an
interleaved and load-labelled matrix to re-run on an idle host, and the rule in decision~008 for what counts as a win (a median
over rounds above the baseline on both tests for the model the flag serves).

\begin{table}[t]
\centering
\caption{Engine flags on the shared host: median tokens/s over three interleaved rounds, relative to the baseline
(\texttt{-t 4}, flash attention auto, micro-batch 512, poll 50). Load before a row: 8--67, median 14.}
\label{tab:p14-flags}
\scriptsize
\begin{tabular}{@{}lrrrr@{}}
\toprule
Flag & S1 pp512 & S1 tg64@1024 & S2 pp512 & S2 tg64@1024 \\
\midrule
baseline (tokens/s) & 87.9 & 8.4 & 35.0 & 3.45 \\
\texttt{-fa off} & 0.79$\times$ & 0.66$\times$ & 0.88$\times$ & 0.25$\times$ \\
\texttt{-fa on} & 0.79$\times$ & 0.31$\times$ & 0.58$\times$ & 0.18$\times$ \\
\texttt{-ub 256} & 0.96$\times$ & 0.69$\times$ & 0.42$\times$ & 0.26$\times$ \\
\texttt{-ub 1024} & 0.96$\times$ & 0.94$\times$ & 0.55$\times$ & 0.32$\times$ \\
\texttt{-t 3} & 0.86$\times$ & 1.02$\times$ & 0.64$\times$ & 0.51$\times$ \\
\texttt{--poll 0} & 0.93$\times$ & 0.56$\times$ & 0.64$\times$ & 0.50$\times$ \\
\texttt{-lm mlock} & 1.18$\times$ & 0.70$\times$ & 0.73$\times$ & 0.35$\times$ \\
\bottomrule
\end{tabular}
\end{table}

\subsection{Co-scheduling the two servers}
\label{sec:p14-cosched}

Both servers run four threads on four cores. \texttt{scripts/cosched\_ab.py} starts a temporary pair per variant and measures
System~2 decoding alone and with System~1 decoding beside it, and the reverse, two interleaved repetitions each
(Table~\ref{tab:p14-cosched}). When both decode, each keeps about half to two thirds of its solo rate and the pair writes
7.1--7.8 tokens/s together in every variant: the cores are the bottleneck, and the flags only move the split. Polling off does
not help and priorities are unavailable on macOS. Freezing the quick model while the thorough one works would cut the pair's
throughput and the other user's time to first token for no gain in the thorough model's, so nothing was built; the serving
script is unchanged.

\begin{table}[t]
\centering
\caption{Co-scheduling on four cores (shared host, load 6--11; medians of two repetitions). ``Keeps'' is the rate with the
other server busy over the rate alone.}
\label{tab:p14-cosched}
\scriptsize
\begin{tabular}{@{}lrrrrrrr@{}}
\toprule
Variant & S2 alone & S2+S1 & keeps & S1 alone & S1+S2 & keeps & both \\
\midrule
base & 3.62 & 2.29 & 63\,\% & 9.87 & 5.21 & 53\,\% & 7.50 \\
\texttt{--poll 0} & 3.67 & 2.37 & 65\,\% & 10.02 & 5.45 & 54\,\% & 7.82 \\
S1 \texttt{-t 2} & 4.33 & 2.02 & 47\,\% & 6.88 & 5.49 & 80\,\% & 7.52 \\
S1 \texttt{-t 3} & 4.11 & 2.00 & 49\,\% & 9.40 & 5.42 & 58\,\% & 7.42 \\
\texttt{--poll 0}, S1 \texttt{-t 2} & 6.37 & 1.69 & 27\,\% & 9.55 & 5.40 & 57\,\% & 7.09 \\
\bottomrule
\end{tabular}
\end{table}

\subsection{The model switch}
\label{sec:p14-switch}

Inside a request the switch is the gate: draft, exact check, rescue batch, decision, System~2. \texttt{scripts/handoff\_latency.py}
times it on the live gateway from the SSE events and both servers' metrics, with the rescue on and off, interleaved.
Four requests whose draft fails its exact check, two with the rescue on and two off, on the shared host (load 7--12): the gate's
own work takes 0.8--0.9\,s (26.9--34.5\,s when the four-sample rescue batch runs, the batch Phase~13 measured to pay for
itself); from the decision to System~2's first reasoning token 1.1--4.7\,s (median 1.5\,s), with 5--27 new prompt tokens against
262--1,148 cached; and System~1 reports no busy slot after the decision in any run. The switch inside a request is already as
fast as the engine allows: there is nothing to freeze (an idle llama-server costs no CPU), and a prefetch of the System~2
prompt during the draft could save at most the 1.1--4.7\,s above at the price of contending with the draft, so none was built.

Between backends, \texttt{POST /v1/system/backend} now warms the target pair before the swap: every constant prefix (the
system prompt per domain on both models, the agent's prompt) is primed on the new servers, then the active pair is exchanged in
one atomic pointer swap. Requests in flight hold the pair they started on and finish there; new requests find the new pair's
caches warm (about 1\,s to the first token against about 9.8\,s cold, decision~005). The response reports the prefixes warmed,
the warm-up time and the requests still running on the previous pair.

\subsection{A cascade in the harness}
\label{sec:p14-cascade}

Codex's agent core offers three things Phase~11 did not take: parallel tool calls, dispatch while the model still streams, and
retries with backoff of a failed sampling request. The first two have nothing to overlap here (one action per reply; under the
grammar the server stops at the end of the action). The third was missing and is now in: a model call that fails for a reason
other than the run's cap is retried up to twice with Codex's backoff (200\,ms, doubling, jitter $0.9$--$1.1$).

What the harness lacked was \Sys{}'s own idea. The agent ran on System~2 from its first step, while System~1 verified 7 of the
same 10 tasks at about 11\,s a step (Phase~11, round~2) and four System~1 candidates in one request cost 1.7 times one
(Phase~13). With \texttt{harness.cascade} the first step goes to System~1 as one request with \texttt{cascade\_samples} choices
(the prompt prefilled once, the choices decoded in one batch, at temperature 0.7 so they differ); each candidate's action is
applied to its own copy of the files and judged by the exact verifier in order; the first that passes ends the run and the rest
are not judged; when none passes, the best-scoring candidate (per-assert credit) becomes the state, its reply and observation
stay in the transcript, and System~2 continues from step two as before. A candidate is offered no \texttt{run} (it is judged by
its files), and a run asked for on System~1 itself gets no cascade. The A/B interleaves the cascade with the Phase~11 default on
the same ten tasks (\texttt{scripts/harness\_ab.py --arms s2,cascade}).
Table~\ref{tab:p14-cascade} gives the result on the shared host (load 5--16): both arms verify all ten tasks with no malformed
step; the cascade takes 1,467\,s against 1,987\,s ($-26\,\%$) and is faster in 7 of 10 pairs (exact sign test $p=0.17$:
indicative at $n=10$). Six of the ten tasks never reach System~2: a four-candidate System~1 request verifies them in 10--101\,s
(median 34\,s against 85\,s for the same tasks on System~2). The other four pay one System~1 step before System~2's; two still
end faster, two slower, within the variance of System~2's own time (the sequential arm took 342\,s and 322\,s on two of them).
The cascade decodes more tokens in total (5,958 against 4,534) but on the cheap model in one batch; what it saves is System~2
steps, 4 against 12. The default is now \texttt{harness.cascade = true} with four candidates; decision~008 says what re-opens it.

\begin{table}[t]
\centering
\caption{The harness cascade against the Phase~11 default, ten tasks interleaved, System~2 as the run's model (shared host).}
\label{tab:p14-cascade}
\scriptsize
\begin{tabular}{@{}lrrrrr@{}}
\toprule
Arm & Verified & Calls & Steps S1 / S2 & Wall & Median \\
\midrule
Phase~11 default & 10/10 & 12 & 0 / 12 & 1,987\,s & 96\,s \\
Cascade (4 candidates) & 10/10 & 14 & 10 / 4 & \textbf{1,467\,s} & 100\,s \\
\bottomrule
\end{tabular}
\end{table}

\subsection{Summary}

Table~\ref{tab:p14} summarises the phase. Three things changed: the cascade (on), retries with backoff (on), and the warm-first
backend switch (on). Three were measured and left alone, with the method and the re-open rule recorded: the engine flags
(inconclusive on a shared host), co-scheduling (nothing to arbitrate on four cores), and the switch inside a request (already
about 1.5\,s; nothing to freeze, no prefetch worth building). Every suite was re-run on the final code (Section~\ref{sec:sweep}
lists them); the counts are in the decision record.

\begin{table}[t]
\centering
\caption{Phase~14 outcomes.}
\label{tab:p14}
\scriptsize
\begin{tabular}{@{}p{3.2cm}p{1.4cm}p{2.2cm}@{}}
\toprule
Change & Default & Evidence \\
\midrule
Harness cascade (FR-097) & on, 4 & $-26\,\%$ wall, 7 of 10 pairs \\
Retries with backoff (FR-098) & on, 2 & unit tests \\
Warm-first backend switch (FR-099) & on & 9.8\,s against 1\,s cold/warm \\
Engine flags & unchanged & inconclusive, shared host \\
Co-scheduling & unchanged & 7.1--7.8 tok/s in every variant \\
Prompt prefetch for System~2 & not built & switch costs 1.1--4.7\,s \\
\bottomrule
\end{tabular}
\end{table}
